%% proposal.tex -> proposal.pdf
%%
%% hangar: a lightweight module system for underwater vehicles, with a
%% recorded learning loop. Research proposal, in IEEE conference format
%% (IEEEtran, two columns, IEEEtran.bst). Converted from proposal.md;
%% refs.bib holds the 30 cited entries, copied from literature/refs.bib,
%% which is not in the repository.
%%
%%   make            proposal.pdf  (tectonic, or latexmk with XeLaTeX)
%%
%% Built with XeLaTeX: UTF-8 source, fonts through fontspec. TeX Gyre
%% Termes, Heros and Cursor stand in for IEEEtran's Times, Helvetica and
%% Courier, which have no Unicode shapes under XeTeX. (The class selects
%% its Times family once while it loads, before fontspec can run; the four
%% "Font shape TU/ptm ... undefined" warnings in the log come from that and
%% are harmless: every glyph in the PDF is from the TeX Gyre fonts.)
\documentclass[conference]{IEEEtran}

\usepackage{fontspec}
\setmainfont{texgyretermes}[Extension=.otf, UprightFont=*-regular,
  BoldFont=*-bold, ItalicFont=*-italic, BoldItalicFont=*-bolditalic]
\setsansfont{texgyreheros}[Extension=.otf, UprightFont=*-regular,
  BoldFont=*-bold, ItalicFont=*-italic, BoldItalicFont=*-bolditalic]
\setmonofont{texgyrecursor}[Extension=.otf, UprightFont=*-regular,
  BoldFont=*-bold, ItalicFont=*-italic, BoldItalicFont=*-bolditalic]
\usepackage{cite}
\usepackage[hyphens]{url}
\usepackage{booktabs}
\usepackage{array}

\newcolumntype{L}[1]{>{\raggedright\arraybackslash}p{#1}}

\begin{document}

\title{hangar: a lightweight module system for underwater vehicles, with a
recorded learning loop}

\author{\IEEEauthorblockN{H. Oğuz Bolat}
\IEEEauthorblockA{Department of Geomatics Engineering\\
Yıldız Technical University\\
İstanbul, Türkiye\\
oguz.bolat@std.yildiz.edu.tr}}

\IEEEspecialpapernotice{(Research proposal)}

\maketitle

% Decide before sending: whether to add a note on the use of AI tools in
% preparing this proposal and the literature survey behind it, and how to
% word it.

\begin{abstract}
hangar is a module system for remotely operated vehicles (ROVs) of the
BlueROV2 class, native to Linux and small enough for the vehicle's Raspberry
Pi. Each function of the vehicle is its own process under one manager;
processes exchange fixed-size messages through shared memory, and no reader
can delay a writer. The learning loop is part of this infrastructure, not an
application on top of it. Every message records, by identifier, the input
messages it used. One dataset definition produces the same model inputs live,
offline and inside the running policy. Weights trained on a topside GPU are
checked by the vehicle before use, run in shadow until an operator promotes
the evaluated model by its hash, and are demoted automatically. Any thruster
command can then be traced, by identifiers alone, to the model that produced
it, the rows that trained it, and the sensor readings and pilot commands
behind those rows. hangar supports known deployment-time learning schemes
(pilot takeovers as in HG-DAgger~\cite{kelly2019hgdagger} and
Sirius~\cite{liu2023sirius}; a residual on a classical controller) and makes
them recorded, checked and reversible on a Pi. I propose to implement it in
C11, compare it with ROS~2 on the same Pi, test its failure behaviour on the
bench, and measure what the loop costs and what its checks catch.
\end{abstract}

\begin{IEEEkeywords}
remotely operated vehicles, module system, shared memory, provenance,
deployment-time learning, Raspberry Pi
\end{IEEEkeywords}

\section{Motivation and problem}

On every dive the pilot shows how the vehicle should be driven. Learning from
this data meets three obstacles.

A recording holds messages and timestamps, not which inputs a module used for
each output, so training pairs are rebuilt by matching times, and features
built offline can differ from what the model meets on the vehicle. Different
code paths for training and serving are the most common cause of this
skew~\cite{breck2019datavalidation}: LeRobot trains on decoded video and
deploys on raw frames~\cite{cadene2026lerobot}; AquaBot trains on JPEG frames
and deploys on uncompressed ones, with a time join that can look forward in
time~\cite{liu2025aquabot}. Message-level provenance exists for ROS, carried
in headers~\cite{witte2021provenance} or reconstructed from
traces~\cite{bedard2023messageflow}, but it has not been used to build
datasets.

Moving a model trained topside onto a running vehicle is, in the systems we
know of, a manual step outside the middleware, with nothing on the vehicle to
check the weights or remove a misbehaving model.

ROS~2 brings a DDS participant per process and quadratic discovery to a
computer that also runs the autopilot and encodes video; its authors measure
2.1 to 4.3~MB per added process and recommend composition into one process,
which gives up isolation~\cite{macenski2023composition}.

\section{Proposed approach}

\textbf{Transport and supervision.} A topic is one shared-memory ring per
output, with one writer and any number of readers; a sequence counter per
slot, read in Lamport's order~\cite{lamport1977concurrent}, detects torn
copies. The ring itself is old; iceoryx2 is a current
example~\cite{iceoryx2}. What differs: readers hold no resources, so a dead or
stuck reader needs no cleanup and delays nobody, and schema, provenance and a
staleness flag live in the ring file. The manager restarts a failed module
alone, with back-off, and moves no data; if it crashes or is re-executed for
an upgrade, the modules keep running and are adopted again from a small
registry. This is crash-only design~\cite{candea2003crashonly} and
one-for-one supervision~\cite{armstrong2003thesis}; the recovery from a
registry is ROS Rescue's, for the ROS~1 master~\cite{kaveti2019rosrescue}.

\textbf{Provenance and datasets.} Each message lists up to four input
messages as (topic, sequence number), filled in by the library at publish
time; an output computed from a stale input is flagged, transitively. This is
recorded usage, not proven derivation~\cite{provdm2013}. A dataset spec names
an anchor topic and the inputs joined to it, by provenance or by a bounded
time join. The same join code runs in the live pairer on the vehicle, in the
offline build, and in the policy that builds its inputs. Live and offline
rows are identical except where a \texttt{late} flag marks a message that
arrived after the live row was emitted. A bridge mirrors topics over the
tether with their sequence numbers, so a provenance link means the same
message on both sides, and pilot commands take priority over dataset rows and
weight files.

\textbf{Policies.} A policy module runs a small inference engine that
allocates all memory at load. The engine is standard engineering; the
handover is not. A weights file carries an architecture hash, probe inputs
with expected outputs and tolerances, and the lineage of its training rows.
The vehicle checks hashes, finite values, probes and step time against its
tick budget, warms the model up beside the running one and swaps between two
ticks; \texttt{revert} restores the previous one. The probe check is
differential testing in CRADLE's sense~\cite{pham2019cradle}, done by the
receiving vehicle rather than offline.

\textbf{Safety.} All control sources pass through one multiplexer (the mux)
that enforces thrust, rate and depth limits per axis; the pilot wins on any
axis. The mux is a Simplex-style switch~\cite{seto1998simplex} with a
hand-written envelope and a latched demotion that only a human
releases~\cite{hobbs2023rta}; SOTER, the closest precedent, switches back
automatically where hangar waits for an operator~\cite{desai2019soter}. A new
model runs in shadow until promoted, and is demoted if its output goes stale
or sits at a limit for too long. One gate per actuator outputs neutral unless
the vehicle is armed and the command is fresh. The envelope bounds commands,
not states; it is not a safety proof. In phase~1 ArduSub still drives the
thrusters, held in MANUAL mode with its failsafes on, and hangar owns the
ground-station link. Version~1 of the loop drives only in residual mode, a
bounded correction added to a classical
controller~\cite{johannink2019residual}, as Neural-Fly did beside an autopilot
on a Pi~4~\cite{oconnell2022neuralfly}.

\section{Contribution and prior work}

A survey of about 110 works preceded this revision. Every part of hangar has
prior art; the claim is the combination and where it is done. To our
knowledge, no system combines the properties above, and none on an embedded
vehicle behind a tether. hangar claims no new learning method, safety method
or loop; Table~\ref{tab:prior} names the closest prior work. The secondary
claims, a smaller footprint than ROS~2 with process isolation kept and safe
failure behaviour, are not new in kind and are reported as measured.

\begin{table*}[t]
\caption{Closest prior work}
\label{tab:prior}
\centering
\footnotesize
\begin{tabular}{@{}L{0.17\textwidth}L{0.79\textwidth}@{}}
\toprule
Work & What it has; what hangar adds, or merely combines \\
\midrule
SERL and HIL-SERL~\cite{luo2024serl,luo2025hilserl};
LeRobot~\cite{cadene2026lerobot} &
The loop's shape is theirs (actor on the robot, learner on a GPU, weights
pushed back while the robot runs, human takeover in the actor); hangar adds
checks on received weights, shadow, promotion, demotion and lineage. \\
Sirius~\cite{liu2023sirius} &
The scheme is Sirius's and HG-DAgger's~\cite{kelly2019hgdagger} (per-sample
``who acted'' labels, retraining in rounds during deployment); hangar records
the per-axis source in the data plane and makes rows, weights and promotions
traceable. \\
Neural Simplex~\cite{phan2020neuralsimplex} &
Shadow mode, retraining while the baseline drives and automatic switch-back,
in simulation; hangar runs shadow on a real vehicle with an operator, not a
test, in the switch. \\
Ford patents~\cite{ford2025onboardvalidation,ford2026invehiclelearning} &
A newly trained model run in shadow on the vehicle and validated against the
previous one; a method, not a measured system, and as read without
hash-bound operator promotion, probe or step-time checks, automatic demotion
or lineage. \\
TFX~\cite{baylor2017tfx}, ML Metadata~\cite{mlmd} &
One transform for training and serving, ``safe to serve'' checks and artefact
lineage, on servers; hangar moves the checks onto the vehicle and extends
lineage from artefacts to messages and actuator commands. \\
Witte and Tichy~\cite{witte2021provenance}; Bédard et
al.~\cite{bedard2023messageflow} &
Message-level provenance in ROS, per field in headers or inferred from
traces; hangar fills fixed-size identifier references at publish time, keeps
them across the tether and builds datasets from them. \\
AquaBot~\cite{liu2025aquabot} &
Pilot demonstrations on a low-cost ROV, behaviour cloning and improvement
from rollouts, with the policy topside on two GPUs; hangar runs the policy on
the vehicle under a gate without AquaBot's train/serve skew, and is not the
first to learn from an ROV pilot. \\
Carlucho et al.~\cite{carlucho2018adaptive} &
Online deep reinforcement learning on a real AUV with full authority to the
network; hangar bounds that authority and records the loop, the supervisor
layer they list as future work. \\
\bottomrule
\end{tabular}
\end{table*}

\section{Evaluation plan}

Table~\ref{tab:measure} lists eight measurements that no surveyed work
reports and that carry the claim. Each is a tool in the repository,
reproducible by a script. Where the design states a budget, that is the pass
criterion; the rest are reported as measured.

\begin{table*}[t]
\caption{Measurements that carry the claim}
\label{tab:measure}
\centering
\footnotesize
\begin{tabular}{@{}L{0.24\textwidth}L{0.72\textwidth}@{}}
\toprule
Quantity & How it is measured \\
\midrule
Train/serve mismatch rate, per dive &
Each policy publishes (anchor sequence, CRC of its input) per tick;
key-joined against the built rows; share of ticks that differ \\
Provenance overhead on a Pi &
Publish latency and memory with 0 and with 4 references, Pi~4 and Pi~5 \\
Loop latency &
Newest row a model trained on to its first promoted output, on one clock \\
Cost of a hot swap &
Tick jitter and missed ticks before, during and after a weights load at
50~Hz \\
What the vehicle-side checks catch &
Probe, tolerance and step-time refusals over mutated weights files (wrong
operator semantics, precision change, corrupted weights) \\
Demotion and envelope-activation rates &
Per dive and per model, from the mux state and the evaluation topic; false
demotions listed \\
Lineage completeness &
Share of thruster commands traceable to messages; forward queries from a bad
message to the models it reached \\
Isolation at what cost &
Memory (PSS) and latency of N isolated modules against composed
ROS~2~\cite{macenski2023composition} and iceoryx2~\cite{iceoryx2}, same
Pi~4, message sizes and rates \\
\bottomrule
\end{tabular}
\end{table*}

Secondary: a fault campaign with propellers off (SIGKILL per module, 100
runs; time to \texttt{ready}; time to neutral at the PWM output; manager
crash and upgrade with calls in flight); pilot overrides per promoted model;
residual adaptation after a mass change, in water; and the electric-eye port
hashed against v0.3.1, with video latency no worse than 3.8~ms at
640$\times$480 and 8.3~ms at 1280$\times$720. The learning tasks run on
recorded data, then in simulation and on the bench, and in water only when a
vehicle and a tank are available.

\section{Work plan (draft, to be adjusted)}

\begin{enumerate}
\item \textbf{Language and literature.} Done; see Status.
\item \textbf{Core.} Topics, manager, calls, parameters; budgets as tests;
  the benchmark on a Pi.
\item \textbf{Vehicle integration.} Safety, mux and gate; ArduSub in
  simulation; tether bridge; electric-eye supervised unchanged.
\item \textbf{Data.} Recorder, pairer, offline build; the
  live-equals-offline test.
\item \textbf{Learning loop.} Inference engine, weights file, lineage; both
  tasks on recorded data.
\item \textbf{Bench and water.} Fault campaign; both tasks on the bench,
  then in water if possible; electric-eye ported and hashed.
\item \textbf{Write-up.}
\end{enumerate}

\section{Risks}

\begin{itemize}
\item \textbf{Scope and hardware.} Fifteen design documents for one student;
  each phase ends in a measurable result, and the loop can be shown in
  simulation and on the bench without wet runs. The work needs a Pi~4 or 5, a
  topside computer with a GPU and, for wet runs, a vehicle and a tank.
\item \textbf{The depth-hold example is wrong as designed.} It learns total
  heave and runs as a residual, so it would sit at its cap. Its label, the
  heave stick command, is not heave thrust: ArduSub's stock mixer scales all
  outputs down against the largest input to avoid
  saturation~\cite{ng2024ardusub}, so the command is not thrust once other
  axes are active. Label and cap are being reworked; the cap becomes a
  quantity measured on the vehicle.
\item \textbf{Narrowed claims.} After the survey and two adversarial reviews,
  ``byte-identical rows'' became ``identical except where flagged late'';
  ``computed from'' became ``used''; pilot imitation is HG-DAgger, not DAgger;
  the envelope gives no state guarantee; and version~1 drives only in
  residual mode. Three related works and the full claims of the Ford patents
  are still unread, so the claim may narrow further.
\item \textbf{Weak policies and footprint.} A few dives may not train a good
  policy, and the survey found no precedent for a small network on a Pi CPU
  imitating a pilot from camera frames at 50~Hz. If hangar is not much
  lighter than composed ROS~2, that is reported.
\end{itemize}

\section{Current status}

The design is complete in fifteen documents (00--14), including deployment
(signed releases, upgrade over the tether) and simulation (two managers on
one machine, a 6-DOF model, a camera fed from recorded footage, scripted
faults), revised after two adversarial reviews and a simplification pass. The
language is decided: C11, static musl binaries built with zig cc. A spike
built the same topic and manager slice in C and in Rust, on x86-64 only: the
C manager is 50~KB static against 435~KB, idle resident memory 136~KB against
480~KB, publish 29~ns against 35~ns; wake latency is 5.4~\textmu{}s at the
median for both, and both showed 0 torn reads over 250 million and 0 lost
wake-ups. Nothing has run on a Pi. electric-eye v0.3.1~\cite{bolat_eeye},
whose frame bus, process handling, discovery and configuration hangar builds
on, is forked into the repository and builds with zig cc natively and as
static AArch64 musl binaries with NEON; its 17 test suites pass natively and
under qemu-aarch64. OpenMP is unavailable in static musl builds, so its
operator loops will use a thread pool when ported. Licences: GPL-3.0-or-later
for the system, LGPL-3.0-or-later for the module library and the Python
package. Nothing has run on a vehicle.

% Balances the columns of the last page (IEEEtran's mechanism: a column
% break before the given reference). Adjust if the text changes.
\IEEEtriggeratref{23}
\bibliographystyle{IEEEtran}
\bibliography{refs}

\end{document}
